\section{Rational Unification}

In this section we introduce equation systems and the unification algorithm, and prove its termination and correctness properties.

For convenience, we use the notation $A^? = \{ \bot \} \cup A$ for some set $A$ and call it ``optional $A$''. Also we use the notation $\pi_i$ for the $i$-th projection of a tuple; for example, $\pi_1 : A \times B \rightarrow A$.


\subsection{Equation System}

We define equation $e \in \E$ as a pair of:
\begin{itemize}
    \item a \textit{pointed} finite set $X \subset \V$ named ``left-hand side'' with the base point named ``representing variable'';
    \item an optional non-variable term $t \in T^? \setminus \V$ named ``right-hand side''.
\end{itemize}

For convenience, we use $\lhs : \E \rightarrow 2^\V$, $\repr : \E \rightarrow \V$ and $\rhs : \E \rightarrow T^?$ for equation components respectively. Also we use the notation $\underline{x}, y, z \equiv t$ in examples where $\lhs(\underline{x}, y, z \equiv t) = \{ x, y, z \}$, $\repr(\underline{x}, y, z \equiv t) = x$ and $\rhs(\underline{x}, y, z \equiv t) = t$.

We define equation system $\xi \in \Xi$\inRocq{wf_eqsys} as a finite set of equations $\Xi \subset 2^\E$ with disjunctive left-hand sides. For convenience we use the notation $e, \xi$ for the equation system extension that formally means $\{ e \} \cup \{ e' \in \xi \mid \lhs(e) \cap \lhs(e') = \emptyset \}$, i.~e. add $e$ in $\xi$ in place of the equations with intersecting left-hand sides.

We consider the given equation definition as a representation of the special case of equation system in conventional meaning. For example:
\begin{equation*}
    \begin{array}{c|c}
        \text{Conventional equation system} & \text{The our equation} \\[2pt]
        \left\{
        \begin{array}{l@{{}={}}l}
            x & f(z) \\
            x & y \\
        \end{array}
        \right. &
        \underline{x}, y \equiv f(z) \\
    \end{array}
\end{equation*}

Under this sense, we may see that trivial equations of the form $\underline{x} \equiv \bot$ doesn't equate anything so we consider them illegal, i.~e. $\forall x \in \V, \xi \in \Xi. ~ \underline{x} \equiv \bot \not\in \xi$. On the other hand, we preserve their existence in $\E$ for convenience in the following definition.

The main operation defined on equation systems is $\walk : \Xi \times \V \rightarrow \E \times T$\inRocq{wf_eqsys_walk}. This operation looks up a variable in an equation system returning a corresponding pair of equation and term. More formally:
\begin{equation*}
    \walk(\xi, x) = \left\{
    \begin{array}{lL}
         (\underline{x} \equiv \bot, x) & if no such $e \in \Xi$ so $x \in lhs(e)$, \\
         (e, repr(e)) & if $e \in \Xi, x \in \lhs(e)$ and $rhs(e) = \bot$, \\
         (e, rhs(e)) & otherwise. \\
    \end{array}
    \right.
\end{equation*}

Basically, we may consider the our definition of equation system as just an enhanced triangular substitution~\cite{baader2001unification} w.~r.~t. \walk. For example:
\begin{equation*}
    \begin{array}{c|c}
        \text{Triangular substitution} & \text{Equation system} \\[2pt]
        \left\{
        \begin{array}{l@{{}\mapsto{}}l}
            x & f(z) \\
            y & x \\
            z & w \\
        \end{array}
        \right. &
        \left\{
        \begin{array}{l@{{}\equiv{}}l}
            \underline{x}, y & f(z) \\
            z, \underline{w} & \bot \\
        \end{array}
        \right. \\
    \end{array}
\end{equation*}

There is a number of trivial lemmas about \walk:
\begin{itemize}
    \item if $\walk(\xi, x) = (e, y)$ and $y \in \V$ then $\repr(e) = y$\inRocq{wf_eqsys_walk_var};
    \item if $\walk(\xi, x) = (e, t)$ and $t \not\in \V$ then $e \in \xi$ and $rhs(e) = t$\inRocq{wf_eqsys_walk_last_nonvar};
    \item $\walk(\xi, x) = \walk(\xi, \repr(\pi_1(\walk(\xi, x))))$\inRocq{wf_eqsys_walk_idemp};
    \item if $\walk(\xi, x_1) = (e_1, t_1)$, $\walk(\xi, x_2) = (e_2, t_2)$ and $\repr(e_1) = \repr(e_2)$ then $t_1 = t_2$\inRocq{wf_eqsys_walk_fst_inj}.
\end{itemize}

Naturally, we may consider equation systems as substitutions (denoted as $\llbracket \cdot \rrbracket : \Xi \rightarrow \Sigma$) that applying recursively to the given term. More formally, we may consider a homomorphic extension of \walk w.~r.~t. the second projection of its result and apply it infinitely:
\begin{equation*}
    \begin{array}{r@{{}={}}l}
        x \llbracket\xi\rrbracket & \pi_2(\walk(\xi, x)) \llbracket\xi\rrbracket \\
        f^n(t_1, ..., t_n) \llbracket\xi\rrbracket & f^n(t_1 \llbracket\xi\rrbracket, ..., t_n \llbracket\xi\rrbracket) \\
    \end{array}
\end{equation*}

Consider an application of equation system $\xi$ to variable $x$. There are two main cases:
\begin{itemize}
    \item if $\pi_2(\walk(\xi, x)) = x' \in \V$ have $x \llbracket\xi\rrbracket = x' \llbracket\xi\rrbracket$ that leads to an infinite recursion but we know that $\walk(\xi, x) = (e, x')$, $\repr(e) = x'$ and $(e, x') = \walk(\xi, x) = \walk(\xi, \repr(\pi_1(\walk(\xi, x))) = \walk(\xi, x')$ so we must consider that $x' \llbracket\xi\rrbracket = x'$;
    \item otherwise, we just descend into the right-hand side term.
\end{itemize}

Therefore, equation system application may be defined as a corecursive function\inRocq[\break]{wf_eqsys_image_hlp}:
\begin{algorithmic}
    \Function{ApplyEqSys}{$\xi \in \Xi$, $t \in T^\infty$}
        \If{$x \gets t \in \V$}
            \State $(\_, t) \gets \walk(\xi, x)$
        \EndIf
        \If{$x' \gets t \in \V$}
            \State \Return $x'$
        \ElsIf{$f^n(t_1, ..., t_n) \gets t$}
            \State \Return $f^n(\Call{ApplyEqSys}{\xi, t_1}, ..., \Call{ApplyEqSys}{\xi, t_n})$
        \EndIf
    \EndFunction
\end{algorithmic}

Under the equation system application we may consider equation systems as a special case of substitutions\inRocq{wf_eqsys_to_subst}.

\begin{rocqtheorem}[Equation system is rational substitution]{wf_eqsys_to_subst_rational}
    Given equation system $\xi \in \Xi$, $\llbracket\xi\rrbracket$ is rational substitution.
\end{rocqtheorem}
\begin{proof}
    Given variable $x$ we may state that any subterm of $x \llbracket\xi\rrbracket$ is either $x$ itself or a some of $\xi$ equation subterms under the application of $\llbracket\xi\rrbracket$. Here we consider equation $e \in \E$ subterms as $\repr(e) \cup \subterms(\rhs(e))$ where $\subterms(\bot) = \emptyset$. Since $\xi$ is finite and $\subterms(t)$ for any $t \in T$ too, the set of the possible subterms is finite too.
\end{proof}

The second key operation of equation systems is $\union : \Xi \times \V \times \V \rightarrow \Xi \times (T \times T)^?$\inRocq{wf_eqsys_union}. Given two variables $x, y$ and equation system $\xi$ this operation makes a new equation system from $\xi$ so that $x, y$ lay in the same equation. More formally:
\begin{algorithmic}[1]
    \Function{UnionEqSys}{$\xi \in \Xi$, $x, y \in \V$}
        \State $(e_x, \_) \gets \walk(\xi, x)$
        \State $(e_y, \_) \gets \walk(\xi, y)$
        \If{$\repr(e_x) = \repr(e_y)$}
            \State \Return $(\xi, \bot)$
            \Comment{already in the one equation or equals}
        \Else
            \If{$\rhs(e_x) = \bot$ \textbf{and} $\rhs(e_y) = \bot$}
                \State \Return $([\lhs(e_x) \cup \lhs(e_y) \equiv \bot, \xi], \bot)$
            \ElsIf{$\rhs(e_x) = \bot$}
                \State \Return $([\lhs(e_x) \cup \lhs(e_y) \equiv \rhs(e_y), \xi], \bot)$
            \ElsIf{$\rhs(e_y) = \bot$}
                \State \Return $([\lhs(e_x) \cup \lhs(e_y) \equiv \rhs(e_x), \xi], \bot)$
            \Else
                \State \Return $([\lhs(e_x) \cup \lhs(e_y) \equiv \rhs(e_x), \xi],(\rhs(e_x), \rhs(e_y)))$
            \EndIf
        \EndIf
    \EndFunction
\end{algorithmic}

When we add a new equation at lines 7-14 we may choose as a representing variable as $\repr(e_x)$ as $\repr(e_y)$. The same is true about the right-hand side of the new equation at line 14: it could be either $\rhs(e_x)$ or $\rhs(e_y)$. As we can see, the second element of the \union result is intended to inform a caller about a collision in right-hand sides.

We consider the following properties of \union and left them without a human-readable proof\inRocq[\break]{wf_eqsys_union_prop}:
\begin{itemize}
    \item if $\union(\xi, x, y) = (\xi', \bot)$ then $\xi'$ is a minimal unifying extension of $\xi$ for $x, y$ (mostly by \autoref{lemma:unbound-unifier});
    \item if $\union(\xi, x, y) = (\xi_1, (t_1, t_2))$ and $\xi_2$ is a minimal unifying extension of $\xi_1$ for $t_1, t_2$ then $\xi_2$ is a minimal unifying extension of $\xi$ for $x, y$ too;
    \item if $\union(\xi, x, y) = (\xi_1, (t_1, t_2))$ and $\xi_2$ is a minimal unifying extension of $\xi$ for $x, y$ then $\xi_2$ is a minimal unifying extension of $\xi_1$ for $t_1, t_2$ too.
\end{itemize}

These properties are required for the unification algorithm correctness and completeness and some other properties required for the termination will be discussed below.


\subsection{Common Part}

We define ``common part'' of two terms $\commonpart : T \times T \rightarrow T^?$ inductively\inRocq{is_common_part}:
\begin{equation*}
    \begin{array}{r@{{}={}}l}
        \commonpart(x, \_) & x \\
        \commonpart(\_, y) & y \\
        \commonpart(f^n(l_1, ..., l_n), f^n(r_1, ..., r_n)) & f^n(\commonpart(l_1, r_1), ..., \commonpart(l_n, r_n)) \\
    \end{array}
\end{equation*}

Given definition is intentionally non-total since two terms may not have a common part (in case of failure the result is $\bot$) and slightly non-deterministic to allow optimizations. Independently of the concrete implementation, we consider \commonpart to be a total deterministic function\inRocq{common_part}.

There is the connection between existence of common part and unifiability:
\begin{itemize}
    \item $\commonpart(t_1, t_2) = \commonpart(t_2, t_1)$\inRocq{is_common_part_sym};
    \item if some substitution $\sigma$ is a unifier of $t_1, t_2 \in T$ and $\commonpart(t_1, t_2) = t$ then $\sigma$ is a unifier of $t_1, t$\inRocq{is_common_part_inf_unifiable};
    \item if some substitution $\sigma$ is a unifier of $t_1, t_2 \in T$ then $\commonpart(t_1, t_2) \neq \bot$\inRocq[\break]{is_common_part_inf_unifiable_inv}.
\end{itemize}


\subsection{The Unification Algorithm}

We define the unification algorithm as a recursive function $\unify : \Xi \times T \times T \rightarrow \Xi^?$\inRocq{rational_unification}:
\begin{algorithmic}[1]
    \Function{RationalUnify}{$\xi \in \Xi$, $t_1, t_2 \in T$}
        \If{$x \gets t_1 \in \V$ \textbf{and} $y \gets t_2 \in \V$}
            \State $(\xi, ts) \gets \union(\xi, x, y)$
            \If{$(t_1, t_2) \gets ts$}
                \State \Return $\Call{RationalUnify}{\xi, t_1, t_2}$
            \Else
                \State \Return $\xi$
            \EndIf
        \ElsIf{$x \gets t_1 \in \V$}
            \State $(e, \_) \gets \walk(\xi, x)$
            \If{$\rhs(e) = \bot$}
                \State \Return $\lhs(e) \equiv t_2, \xi$
            \Else
                \State $t \gets \commonpart(\rhs(e), t_2)$
                \If{$t = \bot$}
                    \State \bf failure
                \Else
                    \State $\xi \gets \lhs(e) \equiv t, \xi$
                    \State \Return \Call{RationalUnify}{$\xi, \rhs(e), t_2$}
                \EndIf
            \EndIf
        \ElsIf{$y \gets t_2 \in \V$}
            \State \Return \Call{RationalUnify}{$\xi, y, t_1$}
        \ElsIf{$f^n(l_1, ..., l_n) \gets t_1$ \textbf{and} $g^m(r_1, ..., r_m) \gets t_2$}
            \If{$f^n \neq g^m$}
                \State \bf failure
            \Else
                \For{$i = 1, ..., n$}
                    \State $\xi \gets \Call{RationalUnify}{\xi, l_i, r_i}$
                \EndFor
                \State \Return $\xi$
            \EndIf
        \EndIf
    \EndFunction
\end{algorithmic}

\begin{rocqtheorem}[Rational unification is correct]{rational_unification_correct}
    Given $\xi \in \Xi$ and $t_1, t_2 \in T$. If $\unify(\xi, t_1, t_2) = \xi'$ and $\xi' \neq \bot$ then $\xi'$ is a minimal unifying extension of $\xi$ for $t_1, t_2$.
\end{rocqtheorem}
\begin{proof}
    Induction by the inference tree using the properties of \union.

    In the case when $t_1 \in \V$ (let's name it $x$) and $t_2 \not\in \V$ we may prove that:
    \begin{itemize}
        \item if $\walk(\xi, x) = (e, \_)$ and $\rhs(e) = \bot$ prove by \autoref{lemma:unbound-unifier};
        \item otherwise, if there is a minimal unifying extension of $\xi'$ for $\rhs(e)$ and $t_2$ then it is a minimal unifying extension of $\xi$ for $x$ and $t_2$ similarly to \union.
    \end{itemize}
\end{proof}

The immediate corollary of this theorem is that $\unify([], t_1, t_2)$ is a MGU of $t_1, t_2$ if not $\bot$\inRocq[\break]{rational_unification_mgu}.

\begin{rocqtheorem}[Rational unification is complete]{rational_unification_complete}
    Given $\xi \in \Xi$ and $t_1, t_2 \in T$. If $\unify(\xi, t_1, t_2) = \bot$ then there is no unifying extension of $\xi$ for $t_1, t_2$.
\end{rocqtheorem}
\begin{proof}
    Induction by the inference tree using properties the of \union.

    In the case when $t_1 \in \V$ (let's name it $x$), $t_2 \not\in \V$, $\walk(\xi, x) = (e, \_)$ and $\rhs(e) \neq \bot$ we may prove that:
    \begin{itemize}
        \item if $\commonpart(\rhs(e), t_2) = \bot$ there is no unifier of $\rhs(e)$ and $t_2$, therefore and a unifying extension of $\xi$ for $x, t_2$;
        \item otherwise, if there is a unifying extension of $\lhs(e) \equiv t, \xi$ (where $t$ is the common part of $t_1, t_2$) for $\rhs(e)$ and $t_2$ then it is a unifying extension of $\xi$ for $x$ and $t_2$ similarly to \union.
    \end{itemize}
\end{proof}

We may also consider the unification of rational terms $t_1, t_2 \in T^R$. Although it isn't straightforwardly comes from the unification of finite terms, it could be obtained using an equation system representation for the given terms:
\begin{enumerate}
    \item Consider equation systems $\xi_1, \xi_2$ and terms $t_1', t_2' \in T$ so that $t_1' \llbracket\xi_1\rrbracket = t_1$ and $t_2' \llbracket\xi_2\rrbracket = t_2$.
    \item Ensure that any introduced fresh variable in $\xi_1$ that isn't presented in $t_1$ isn't presented in $\xi_2$ too and vice-versa. More formally:
    \begin{equation*}
        \left\{
        \begin{array}{l@{{}\implies{}}l}
            \forall e_1 \in \xi_1, x \in \lhs(e_1). ~ x \not\in \subterms(t_1) & \forall e_2 \in \xi_2. ~ x \not\in \lhs(e_2) \\
            \forall e_2 \in \xi_2, x \in \lhs(e_2). ~ x \not\in \subterms(t_2) & \forall e_1 \in \xi_1. ~ x \not\in \lhs(e_1) \\
        \end{array}
        \right.
    \end{equation*}
    \item Iteratively extend $\xi_1$ with $\xi_2$ equations getting the new equation system $\xi$:
    \begin{algorithmic}
        \State $\xi \gets \xi_1$
        \For{$e \in \xi_2$}
            \State $\xi \gets e, \xi$
        \EndFor
    \end{algorithmic}
    \item Finally, perform $\unify(\xi, t_1', t_2')$.
\end{enumerate}

% TODO: proof that?

In practice, it isn't much helpful to unify rational terms directly since we may always represent them with pairs of an equation system and a finite term.


\subsection{Termination of the Algorithm}

To prove the termination of the algorithm we need to introduce some new notions.


\subsubsection{Roots}

We say that a variable $x$ is a ``root'' of an equation system $\xi$ iff it's either:
\begin{itemize}
    \item a representing variable of some $\xi$ equation;
    \item or a variable that is not present in a left-hand side of any $\xi$ equation.
\end{itemize}

More formally, we say that $x$ is a root of $\xi$ iff $\repr(\pi_1(\walk(\xi, x))) = x$.

We define roots of an equation system modulo a finite (and decidable) variable set $X \subset \V$ $\roots_X : \Xi \rightarrow 2^\V$\inRocq{eqsys_roots} as:
\[\roots_X(\xi) = \{ x \in X \mid \repr(\pi_1(\walk(\xi, x))) = x \}.\]

The motivation for the introduction of roots as a separated term (instead of the representing variables) is the unification of unbound variables. For example, when we unify $x$ with $f(y)$ and there is no an equation with $x$ in the left-hand side we are required to add the new equation $\underline{x} \equiv f(y)$. It increases the number of representing variables but doesn't change the number of roots modulo $X$ if $x \in X$.

Naturally, we may consider an equation system as like as there is an equation for every $x \in \V$ just populating it by trivial equations of the form $\underline{x} \equiv \bot$. We may see that the \walk behavior stay unaffected since it's the natural sense of equation systems as like as for substitutions. So a root of $\xi$ is just a representing variable of so populated variant of $\xi$. The issue of this approach is that the cardinality of the representing variable set turns equivalent to the $\N$ cardinality and therefore cannot be properly reduced by one.


\subsubsection{Non-variable size}

We define non-variable size of a term $t \in T$ as a number of its constructors. More formally, we consider the following recursive definition $\size : T \rightarrow \N$\inRocq{term_nonvar_size}:
\begin{equation*}
    \begin{array}{r@{{}={}}l}
        \size(x) & 0 \\
        \size(f^n(t_1, ..., t_n)) & 1 + \sum_{i = 1}^n \size(t_i) \\
    \end{array}
\end{equation*}

We also consider the non-variable size of an equation system $\xi$ as a sum of non-variable sizes of its right-hand sides $\size : \Xi \rightarrow \N$\inRocq{eqsys_nonvar_size}:
\[\size(\xi) = \sum_{\substack{e \in \xi \\ \rhs(e) \neq \bot}} \size(\rhs(e)).\]


\subsubsection{Income}

We define ``income'' of two terms $\income : T \times T \rightarrow \N$ as a recursively defined function\inRocq{rational_unification_income}:
\begin{equation*}
    \begin{array}{r@{{}={}}l}
        \income(x, t_2) & \size(t_2) \\
        \income(t_1, y) & \size(t_1) \\
        \income(f^n(l_1, ..., l_n), g^m(r_1, ..., r_m)) & 1 + \sum_{i = 1}^{\min(n, m)} \income(l_i, r_i) \\
    \end{array}
\end{equation*}

Actually, independently from the actual definition, we only need that:
\begin{itemize}
    \item $\income(x, t) = \income(t, x) = \size(t)$ --- obviously by the definition;
    \item if $\commonpart(t_1, t_2) = t$ then $\income(t_1, t_2) + \size(t) = \size(t_1) + \size(t_2)$ \\
    \inRocq{rational_unification_income_common_part}. In other words, that $\income(t_1, t_2)$ is the sum of the $t_1, t_2$ non-variable sizes without the non-variable size of its' common part.
\end{itemize}


\subsubsection{Unification Task}

We define unification task $\tau \in \T$\inRocq{rational_unification_task} as a triple of an equation system and two terms: $\T = \Xi \times T \times T$.

Now we able to define ``recursive call'' relation between unification tasks $\tau_1, \tau_2 \in \T$ modulo finite (and decidable) $X \subset \V$ denoted as $\tau_1 \reccall \tau_2$\inRocq{rational_unification_recursive_call} as presented on \autoref{fig:recursive-call}. The definition of $\reccall$ is aligned with all 4 possible recursive call kinds in the definition of \unify. More precisely (the item numbers are the listing line numbers):
\begin{enumerate}
    \setitemnumber{5}\item \nameref{rule:union-call};
    \setitemnumber{18}\item \nameref{rule:common-part-call};
    \setitemnumber{20}\item \nameref{rule:swap-var-term-call};
    \setitemnumber{26}\item \nameref{rule:subterm-call}.
\end{enumerate}

\begin{figure}
    \centering
    \begin{subfigure}{.99\linewidth}
        \begin{prooftree}
            \AxiomC{$x, y \in \V$ \quad $|\roots_X(\xi_2)| < |\roots_X(\xi_1)|$}
            \RuleName{UnionCall}
            \UnaryInfC{$(\xi_1, x, y) \reccall (\xi_2, t_1, t_2)$}
        \end{prooftree}
        \RuleCaption\label{rule:union-call}
    \end{subfigure}
    \begin{subfigure}{.99\linewidth}
        \begin{prooftree}
            \AxiomC{$x \in \V$ \quad $t_x, t_y \not\in \V$ \quad $\commonpart(t_x, t_y) = t$}
            \noLine
            \UnaryInfC{$|\roots_X(\xi_1)| = |\roots_X(\xi_2)|$ \quad $\size(\xi_1) + \size(t) = \size(\xi_2) + \size(t_x)$}
            \RuleName{CommonPartCall}
            \UnaryInfC{$(\xi_1, x, t_y) \reccall (\xi_2, t_x, t_y)$}
        \end{prooftree}
        \RuleCaption\label{rule:common-part-call}
    \end{subfigure}
    \begin{subfigure}{.99\linewidth}
        \begin{prooftree}
            \AxiomC{$y \in \V$ \quad $t \not\in \V$}
            \RuleName{SwapVarTermCall}
            \UnaryInfC{$(\xi, t, y) \reccall (\xi, y, t)$}
        \end{prooftree}
        \RuleCaption\label{rule:swap-var-term-call}
    \end{subfigure}
    \begin{subfigure}{.99\linewidth}
        \begin{prooftree}
            \AxiomC{$|\roots_X(\xi_2)| \le |\roots_X(\xi_1)|$}
            \noLine
            \UnaryInfC{$\size(\xi_2) \le \size(\xi_1) + \sum_{j = 1}^{i - 1} \income(l_j, r_j)$}
            \RuleName{SubtermCall}
            \UnaryInfC{$(\xi_1, f^n(l_1, ..., l_n), g^n(r_1, ..., r_n)) \reccall (\xi_2, l_i, r_i)$}
        \end{prooftree}
        \RuleCaption\label{rule:subterm-call}
    \end{subfigure}
    \caption{The definition of $\tau_1 \reccall \tau_2$}
    \label{fig:recursive-call}
\end{figure}

We define a unification task size modulo a finite (and decidable) $X \subset \V$ $\size_X : \T \rightarrow \N \times \N$\inRocq{rational_unification_size} as:
\[\size_X(\xi, t_1, t_2) = (|\roots_X(\xi)|, \size(\xi) + \income(t_1, t_2)).\]

Also we define an equation system unification size modulo a finite (and decidable) $X \subset \V$ $\size_X : \Xi \rightarrow \N \times \N$\inRocq{rational_unification_system_size} as:
\[\size_X(\xi) = (|\roots_X(\xi)|, \size(\xi)).\]

Further we imply lexicographical order on pairs. Under these definitions we may see that every considered recursive call strictly reduces a unification task size (e.~g. it's clear for \nameref{rule:union-call}) or at least doesn't increases it (as in \nameref{rule:swap-var-term-call}). Moreover, we may prove that $\reccall$ is a well-founded relation for any $X \subset \V$.

\begin{rocqlemma}[$\reccall$ is well-founded]{rational_unification_recursive_call_well_founded}
    Given a finite (and decidable) $X \subset \V$, $\reccall$ is a well-founded relation.
    \label{lemma:reccall-wf}
\end{rocqlemma}
\begin{proof}
    Given $(\xi_1, t_1, t_2) \in \T$, we need to show that any recursive call chain is finite:
    \[(\xi_1, t_1, t_2) \reccall \tau_2 \reccall ... \reccall \tau_n.\]
    
    Well-founded induction by the unification task size modulo $X$. In other words, we show that any recursive call chain from $(\xi_1, t_1, t_2)$ is finite if it's true for any $\tau' \in \T$ so that $\size_X(\tau') < \size_X(\xi_1, t_1, t_2)$ (we call it induction hypothesis, IH). Case analysis on $t_1$:
    \begin{itemize}
        \item If $t_1 \in \V$ (let's call it $x$), the only allowed recursive calls are \nameref{rule:union-call} and \nameref{rule:common-part-call}:
        \begin{itemize}
            \item Have $t_2 = y \in \V$, $(\xi_1, x, y) \reccall (\xi_2, t_1', t_2')$ and $|\roots_X(\xi_2)| < |\roots_X(\xi_1)|$. Obviously, we may use IH since the latter implies $\size_X(\xi_2, t_1', t_2') < \size_X(\xi_1, x, y)$ by the lexicographic ordering.
            \item Have $(\xi_1, x, t_2) \reccall (\xi_2, t_x, t_y)$, $t_x, t_y \not\in \V$, $\commonpart(t_x, t_y) = t$, $|\roots_X(\xi_1)| = |\roots_X(\xi_2)|$ and $\size(\xi_1) + \size(t) = \size(\xi_2) + \size(t_x)$.
            \begin{enumerate}
                \item Let's show that $\size_X(\xi_1, x, t_y) = \size_X(\xi_2, t_x, t_y)$:
                \begin{itemize}
                    \item Since $|\roots_X(\xi_1)| = |\roots_X(\xi_2)|$, we only need to show that $\size(\xi_1) + \income(x, t_y) = \size(\xi_2) + \income(t_x, t_y)$.
                    \item By the definition of \income, $\income(x, t_y) = \size(t_y)$.
                    \item By the property of common part and \income, $\income(t_x, t_y) + \size(t) = \size(t_x) + \size(t_y)$.
                \end{itemize}
                \begin{equation*}
                    \begin{array}{r@{{}={}}l}
                        \size(\xi_1) + \underline{\income(x, t_y)} + \size(t) & \size(\xi_1) + \underline{\size(t_y) + \size(t)} \\
                        & \underline{\size(\xi_1) + \size(t)} + \size(t_y) \\
                        & \size(\xi_2) + \underline{\size(t_x) + \size(t_y)} \\
                        & \size(\xi_2) + \income(t_x, t_y) + \size(t).
                    \end{array}
                \end{equation*}
                \item Let's consider the possible recursive calls from $(\xi_2, t_x, t_y)$. Since $t_x, t_y \not\in \V$, it's only \nameref{rule:subterm-call}. We may see that any \nameref{rule:subterm-call} recursive call strictly reduces the terms \income because of the top-level constructors so we may use IH.
            \end{enumerate}
        \end{itemize}
        \item If $t_1 \not\in \V$, the only allowed recursive calls are \nameref{rule:swap-var-term-call} and \nameref{rule:subterm-call}. The first one just swaps terms of the task and immediately turns into the previous case. The latter always strictly reduces the terms \income because of the top-level constructors so we may use IH.
    \end{itemize}
\end{proof}


\subsubsection{Proof}

Finally, we define variables of an equation system $\vars : \Xi \rightarrow 2^\V$\inRocq{eqsys_vars} as the set of variables presented anywhere in the equation system.

We may prove the following properties of \union\inRocq{wf_eqsys_union_prop}:
\begin{itemize}
    \item if $\union(\xi, x, y) = (\xi', \bot)$, $\vars(\xi) \subseteq X$ and $x, y \in X$ then $|\roots_X(\xi')| \le |\roots_X(\xi)|$;
    \item if $\union(\xi, x, y) = (\xi', \bot)$ then $\size(\xi') \le \size(\xi)$;
    \item if $\union(\xi, x, y) = (\xi', \bot)$ then $\vars(\xi') \subseteq \{ x, y \} \cup \vars(\xi)$;
    \item if $\union(\xi, x, y) = (\xi', (t_1, t_2))$ and $\vars(\xi) \subseteq X$ then $|\roots_X(\xi')| = |\roots_X(\xi)| + 1$;
    \item if $\union(\xi, x, y) = (\xi', (t_1, t_2))$ then $\vars(\xi') \subseteq \vars(\xi)$ and there are $e_1, e_2 \in \xi$ so $\rhs(e_1) = t_1$ and $\rhs(e_2) = t_2$.
\end{itemize}

Using these properties we firstly prove that if $\unify(\xi, t_1, t_2) = \xi'$ and $\xi' \neq \bot$ then $\vars(\xi') \subseteq \vars(t_1) \cup \vars(t_2) \cup \vars(\xi)$\inRocq{rational_unification_result_vars} where $\vars : T \rightarrow 2^\V$ \space\inRocq{fv_term} is defined similarly to the \vars of an equation system.

\begin{rocqtheorem}[Rational unification terminates]{rational_unification_exists}
    Given $\xi \in \Xi$ and $t_1, t_2 \in T$. The computation of $\unify(\xi, t_1, t_2)$ terminates.
\end{rocqtheorem}
\begin{proof}
    We need to prove the more general fact here: for any finite (and decidable) $X \subset \V$ so that $\vars(\xi) \cup \vars(t_1) \cup \vars(t_2) \subseteq X$ the algorithm terminates and if its result $\xi' \neq \bot$, $\size_X(\xi') \le \size_X(\xi, t_1, t_2)$. It means that we able to increase the equation system size up to $\income(t_1, t_2)$.
    \begin{itemize}
        \item This fact straightforwardly comes from the properties of \union and the previous lemma, by well-founded induction on $\reccall$.
        \item Finally, we instantiate the generalized fact with $X = \vars(\xi) \cup \vars(t_1) \cup \vars(t_2)$.
    \end{itemize}
\end{proof}
